\section{Related Work}
Agent Skills combine instructions with executable resources, so their security
depends on both source content and the current task. Liu et al. document
confirmed malicious Skills~\cite{liu}, while repository-aware analysis warns
against treating suspicious content as confirmed maliciousness~\cite{context}.
Skill-Inject examines attacks inside Skill files~\cite{skillinject};
InjecAgent studies indirect injection through tool outputs~\cite{injecagent}.
SkillMutator extends this threat to interactions between language and
code~\cite{skillmutator}. These studies motivate inspecting the complete package
and retaining the task context, but their findings do not establish attack
prevalence in our constructed cases.

AgentDojo and Agent Security Bench evaluate harmful actions alongside normal
task completion~\cite{agentdojo,asb}. We follow this evaluation principle by
including permitted communication and deletion as well as attacks. Blocking
every sensitive action therefore cannot count as a useful defence. Our finite
synthetic environment tests specific mechanisms; it is not a reproduction of
either benchmark or a measurement of real Skill ecosystems.

StruQ separates trusted instructions from untrusted data through structured
inputs and training~\cite{struq}. IPIGuard uses tool dependencies to constrain
execution~\cite{ipiguard}, and DRIFT combines task-based planning with dynamic
validation~\cite{drift}. AttriGuard tests whether tool calls survive reduced
external control through counterfactual replay~\cite{attriguard}. Such evidence
addresses why an action occurs; causal support alone does not establish user
permission. Our graph presentation ablation does not implement IPIGuard.

SkillScope is the closest Skill-specific starting point~\cite{skillscope}. It
links instruction and code actions, checks authorization and necessity under
concrete tasks, and constrains validated excess actions. We adapt this framing
while separating necessity in the original implementation from necessity within
a finite alternative catalogue. An unsafe prerequisite can prevent a modified
program from finishing without becoming authorized. We do not claim to use
verified official SkillScope code or reproduce its full system.

Progent already generates task-based symbolic policies and checks policy
updates with an SMT solver~\cite{progent}. CaMeL separates control from untrusted
data and tracks information-flow restrictions~\cite{camel}. VIGIL enforces
behavioural specifications over execution traces~\cite{vigil}, while shielding
provides an earlier foundation for external action enforcement~\cite{shield}.
These precedents rule out a claim to the first task-aware policy or runtime
guard. Our contribution is a bounded combination and controlled evaluation:
source-linked actions, supplied task permissions, finite flow rules, executable
necessity checks, and recovery under unchanged authority and quota. Equal
information and simulation access let us test whether recovery gains come from
the selector itself. This scope does not imply arbitrary-script safety or
general superiority over the prior systems.
